\Qhead{6.5}{Основные задачи нелинейной фильтрации и синтеза систем}{Основные задали нелинейной фильтрации и синтеза систем.}
\Subfile{6.5.1}{sections/06_detection/q05/s01}

\recap{Нелинейная фильтрация распространяет и обновляет всю условную плотность. Линейно-гауссовский случай замыкается на среднем и ковариации; в общем случае нужны приближения и контроль грубых ошибок.}
\board{Прогноз плотности, байесовское обновление и инновационный член; отличие точного фильтра от EKF.}
\oral{Почему одно гауссовское приближение может быть плохим при двух возможных фазах сигнала?}
\bait{Можно оставить сохранение нормировки: среднее $h-\hat h$ по апостериорной плотности равно нулю.}
\sources{Сёрккя и Солин~\cite{sarkka}, §§10.3--10.7, особенно печатные с.~206--210; постановка приёма — вопрос~\Qref{6.1}. Условия Ито и независимости шумов указаны явно.}
